\documentclass[10pt,a4paper]{amsart}
\usepackage[margin=19mm]{geometry}
\usepackage{amsmath,amssymb,mathtools}
\usepackage[scr=rsfs,cal=euler]{mathalfa}
\usepackage{xfrac}
\usepackage[unicode,hidelinks,bookmarks=false]{hyperref}
\newcommand{\ie}{i.e.\ }
\newcommand{\cf}{cf.\ }
\newcommand{\resp}{respectively\ }
\newcommand{\Subsection}{\subsection}
\providecommand{\nobreakdash}{\nobreak\hbox{-}}
\setlength{\emergencystretch}{2em}
\multlinegap0pt
\usepackage{bm}
\usepackage{bbm}
\usepackage{dsfont}
\usepackage{xspace}
\usepackage{cancel}
\usepackage{mathtools}
\usepackage{amsthm}
\makeatletter
\@ifundefined{thm}{\newtheorem{thm}{Thm}}{}
\makeatother
\DeclareMathOperator{\Gr}{Gr}
\DeclareMathOperator{\relint}{relint}
\DeclareMathOperator{\sign}{sign}
\DeclareMathOperator{\spanv}{span}
\newcommand{\CPk}{\mathrm{C}_P^k}
\newcommand{\Pkmax}{P^{[k]}_{\max}}
\newcommand{\mF}{\mathcal{F}}
\newcommand{\mS}{\mathcal{S}}
\title[How to stab a polytope]{How to stab a polytope}
\author{Sebastian Seemann, Francesca Zaffalon}
\date{Source: arXiv:2412.00551v1 (CC BY 4.0)}
\begin{document}
\maketitle

\noindent\emph{Source and licence.} How to stab a polytope. Version arXiv:2412.00551v1 (CC BY 4.0), distributed under the Creative Commons Attribution 4.0 International licence, as recorded in the arXiv licence metadata for this exact version and shown on the abstract page https://arxiv.org/abs/2412.00551v1. Full rights notice: licenses/arxiv-2412.00551-CC-BY-4.0.txt.\par\smallskip

\noindent\textit{Editorial scope.} Counted unit: the Thm whose source label is \texttt{thm : sign var} in the article (source0.tex lines 308--311, source label \texttt{thm : sign var}), delivered together with the article's own complete original proof of it (source0.tex lines 312--343, one \texttt{proof} environment of 32 lines). No mathematics is written, repaired or re-derived: statement and proof are the article's own text line for line. Required context retained uncounted: none. Every definition, display and label of those lines is retained unchanged. Omitted from this package and not redistributed: 1--307, 344--742. Nothing inside a retained excerpt is omitted or altered: every retained body is delivered exactly as the article's lines are written, so no omission receipt of any kind accompanies this package. Citations: the retained excerpts carry no citation command at all, so no bibliography entry of the article is transcribed and no reference is redistributed. Reference adaptation: none was needed and none was made; every reference inside the delivered unit resolves inside it, so no unresolved reference or citation is produced. Printed slips kept literal: no printed word, symbol, hypothesis or number of the counted statement or of its proof was altered. Layout and class: the article's own class is \texttt{ourlematema} with packages including xcolor, cleveref, mathtools, color, graphicx, soul; the wrapper is the lane's pinned \texttt{amsart} editorial wrapper at 10pt on A4, which restates the theorem-like environments the retained text needs and adds the article's own macro definitions that the retained text needs (\texttt{\textbackslash CPk}, \texttt{\textbackslash Gr}, \texttt{\textbackslash Pkmax}, \texttt{\textbackslash mF}, \texttt{\textbackslash mS}, \texttt{\textbackslash relint}, \texttt{\textbackslash sign}, \texttt{\textbackslash spanv}), with extra wrapper packages (bm, bbm, dsfont, xspace, cancel, mathtools). The delivered numbering is the unit's own. Assets: no retained span contains a figure environment, a graphics-inclusion command, a PDF-page inclusion, a table or a TikZ picture; no image file is copied into this package, the package declares no asset dependency and no image file is redistributed with it. Licence: version arXiv:2412.00551v1 of this article is distributed under the Creative Commons Attribution 4.0 International licence, as recorded in the arXiv licence metadata for this exact version and shown on the abstract page https://arxiv.org/abs/2412.00551v1. A full rights notice is delivered at licenses/arxiv-2412.00551-CC-BY-4.0.txt. Characters: every retained excerpt is pure ASCII, so the delivered engine, XeLaTeX with its default Latin Modern text font, reports no missing character.
\begin{thm}\label{thm : sign var}
    The following holds:
    \[ \mS_P(V) = \{ W\in \Gr(k,n) \mid \sign(\CPk|_{\mF(V)}(V)) \equiv \sign(\CPk|_{\mF(V)}(W)) \}. \]
\end{thm}

\begin{proof}
    Let $V\in \Pkmax$ with determining faces of the corresponding chamber given by $\mF_P(V) = \{F_1,\ldots,F_r\}$. 
    Let $W\in \mS_P(V)$, we want to prove that the vector of Chow forms evaluated at $W$ has the same sign vector as the one of $V$ when restricting to the facet of the determining faces.

    The intersection $V\cap P$ is a $(k-1)$-dimensional polytope given as the convex hull of $V\cap F_i = v_i$, for $i=1,\ldots,r\geq k$. Then $V$ can be represented as a $k\times n$ matrix containing $k$ of the vectors $v_i$ as its rows. Without loss of generality, we can assume that these are the first $k$ points $v_i$. We will denote this matrix by $V$. We want to show that we can move the points $v_i$ to the points of intersection $W\cap F_i$ without changing the signs of the Chow forms.

    Let $w_i=W\cap F_i$ and consider the matrix
    \[ U^1_t:=\begin{pmatrix}
    (1-t)v_1 + tw_1 \\
    v_2\\
    \vdots   \\
    v_{k} \\
    \end{pmatrix}. \]
    For every $t\in [0,1]$, the space defined by the matrix $U^1_t$ lies in $\Gr(k,n)$, by convexity of the polytope.
    Hence the map $t \mapsto U^1_t$ is a well-defined path in $\Gr(k,n)$. Moreover, $(1-t)v_1 + tw_1 \in \relint(F_1)$ for each $t\in [0,1]$. The Pl\"ucker coordinates are linear in $t$ along this path. Hence, each Chow form of the facets of $F_1$ is non-zero on each $U_t$, for $t\in[0,1]$. Since the Chow forms are continuous and linear in the Pl\"ucker coordinates, this implies that their sign does not change in the path from $V=U^1_0$ to $U^1_1$. Repeat the same operation for $i=2,\ldots,k$ by defining 
    \[ U^i_t:=\begin{pmatrix}
    w_1 \\
    \vdots\\
    w_{i-1}\\
    (1-t)v_i + tw_i\\
    v_{i+1}\\
    \vdots   \\
    v_{k} \\
    \end{pmatrix}. \]
    In each path we never cross any boundaries of $F_i$, hence the sign of the Chow forms are fixed for every facet of $F_i$. By repeating the same operation for any $k$-subset of the determining faces $F_1,\ldots, F_r$ we obtain the desired result, i.e.
    \[ \sign(\CPk|_{\mF(V)}(V)) = \sign(\CPk|_{\mF(V)}(W)). \]

    \vspace{5pt}
    Suppose now that the sign characterization is true for some $W\in \Gr(k,n)$ and assume by contradiction that $W\not\in \mS_P(V)$. Then there exists some face $F_i\in \mF_P(V)\setminus \mF_P(W)$. Without loss of generality, we can assume that $i=1$. Let $w_i = W\cap \spanv(F_i)$ and define the matrices $U_t^1$ as above. 
    We can choose representative of $w_i$ such that the signs of $C_P(V)|_{\mF(V)}$ and $C_P(W)|_{\mF(V)}$ are the same in affine space.
    By assumption, the intersection of $U_t^1$ with $\spanv(F_1)$ moves outside of $F_1$, hence there exists a $t_0\in [0,1]$ such that $U^1_{t_0}$ intersects a facet $G$ of $F_1$. Since the Pl\"ucker coordinates of $U_t$ are linear in $t$ and the Chow forms are linear in the Pl\"ucker coordinates of $\Gr(k,n)$, $C_G(U^1_t)$ changes sign along the path for $t\in [0,1]$. Since we defined $w_1$ to be the intersection of $W$ with $\spanv(F_1)$, the sign of $C_G$ will stay fixed as we perform the paths defined by $U^2_t, \ldots, U^k_t$. Hence we obtain that $\sign(C_G(W))$ is different (in affine space) from $\sign(C_G(V))$, against the hypothesis. It follows that $W\in \mS_P(V)$.
\end{proof}

\end{document}
